% Research paper shape: IMRaD with a structured abstract.
% Same sections as template.md; the comment above each section is its guidance.
% Standard article class only; compiles with pdflatex (run it twice for references).
\documentclass[11pt]{article}

% Title that states the main finding.
\title{Title that states the main finding}
\author{First Author\thanks{Affiliation. Contact address.}
  \and Second Author\thanks{Affiliation.}}
\date{}

\begin{document}
\maketitle

% Structured abstract: one or two sentences per label, numbers in Results.
\begin{abstract}
\noindent\textbf{Background.} The problem and the gap, in one or two sentences.
\textbf{Methods.} Design, sample and main measure.
\textbf{Results.} The main result with its number and interval.
\textbf{Conclusions.} What the result means, and its main limit.
\end{abstract}

\noindent\textbf{Keywords:} three to six terms a searcher would type.

% Optional: three to five bullets of at most 85 characters, only if the venue asks.
\paragraph{Highlights (optional).}
\begin{itemize}
  \item First finding, in at most 85 characters.
  \item Second finding, in at most 85 characters.
  \item Third finding, in at most 85 characters.
\end{itemize}

% Territory, gap, this paper, contributions, then a one-line map of the sections.
\section{Introduction}
Establish the territory and cite the closest prior work~\cite{placeholder}.
Show the gap. Say what this paper does.
The contributions are:
\begin{enumerate}
  \item first contribution;
  \item second contribution.
\end{enumerate}
Section~\ref{sec:methods} gives the methods, Section~\ref{sec:results} the
results and Section~\ref{sec:discussion} their meaning.

\section{Methods}\label{sec:methods}

% If preregistered, link the protocol and list every departure with its reason;
% if not, say that the study is exploratory.
\subsection{Plan and deviations}
Protocol: link. Deviations: none, or each one with its reason.

% Who or what was studied, how it was chosen, with what; enough to repeat it.
\subsection{Design and materials}
Design, sample, setting and materials.

% Each outcome, how it was measured, and the analysis planned for it.
\subsection{Measures and analysis}
Outcomes, instruments and analysis.

\section{Results}\label{sec:results}

% Headline sentence first, then the number and its interval, then the figure.
\subsection{First claim}
The headline sentence of the first result, with its number and interval
(Figure~\ref{fig:first}).

\begin{figure}[ht]
  \centering
  \fbox{\parbox{0.8\linewidth}{\centering\vspace{3em}Figure placeholder\vspace{3em}}}
  \caption{A caption that says what the figure shows and what to conclude,
    readable without the text.}
  \label{fig:first}
\end{figure}

% One subsection per claim, in the order of the contribution list.
\subsection{Second claim}
The headline sentence of the second result.

\section{Discussion}\label{sec:discussion}

% The main result restated in one paragraph, without new data.
\subsection{Principal finding}
The main result, restated.

% Where the result agrees or disagrees with earlier studies, and why.
\subsection{Comparison with prior work}
Agreement and disagreement with earlier studies.

% What the design cannot show, and how large each threat is.
\subsection{Limitations}
Each limitation and its likely size.

% What follows for practice or for the next study, within what the data support.
\subsection{Implications}
Consequences for practice and for the next study.

% Two or three sentences a reader could quote; no claim absent from the results.
\section{Conclusion}
The conclusion in two or three sentences.

% Where the data, code and environment are, under which licence, how to re-run.
\section*{Data and code availability}
Repository, licence and the command that reproduces the results.

% Name the reporting guideline followed and attach the completed checklist.
\section*{Reporting checklist}
The reporting guideline followed, such as CONSORT or PRISMA, with its checklist attached.

% Funding, conflicts of interest, author contributions, acknowledgements.
\section*{Declarations}
Funding. Conflicts of interest. Author contributions. Acknowledgements.

% In the venue's style; every entry with a DOI, an arXiv id or a stable URL.
\begin{thebibliography}{9}
\bibitem{placeholder} Author, A. (Year). Title of the work. Venue.
  doi:10.0000/placeholder.
\end{thebibliography}

% Material a specialist needs and a general reader can skip.
\appendix
\section{Supplementary material}
Supplementary tables and extra analyses.

\end{document}
